\section{Factorial saturation of the positive-form ideal}

\paragraph{Archive scope.}
This tests a repair for cote 136, physical pp.~117, 119 and 121,
archive footers 116, 118 and 120, supplied scans 117, 119 and 121.
The standard integral model and its Poincaré contraction are the
inherited AQ-136-DP-MODEL. The failure below concerns only (136.PS1).

\paragraph{ARCH-136-POSITIVE-IDEAL-SATURATION.}
Let \(A\) be a torsion-free nonnegative differential PD algebra over
\(S=\Gamma_{\mathbb Z}(T)\), and let \(J=A^{>0}\) be its positive-form ideal.
Inside \(A_{\mathbb Q}\), define
\[
 \operatorname{Sat}_{+}(A)
   =A+\sum_{n\geq2}\frac{1}{n!}J^n .                \tag{136.PS1}
\]
This is a functorial differential PD \(S\)-subalgebra. In form degree \(q\),
only the terms \(n\leq q\) contribute, so the construction uses finitely
many denominators in each fixed form degree. It leaves forms of degrees
zero and one unchanged.

Nevertheless it does not preserve vertical admissibility. For the
standard integral simplex theory \(B\), in original simplex degree two,
\[
 H^2\bigl(\operatorname{Sat}_{+}(B)_2^{w,*}\bigr)
       \cong (\mathbb Z/2)^{\binom w2}
       \qquad(w\geq2).                              \tag{136.PS2}
\]
Thus the failure occurs in every weight \(w\geq2\) in one fixed form degree.

\begin{proof}
The positive-form ideal is stable under the differential. Its powers
therefore are also stable. Products of the summands in (136.PS1) stay in
the indicated sum because
\[
 \frac1{r!}J^r\cdot\frac1{s!}J^s
     \subseteq \binom{r+s}{r}\frac1{(r+s)!}J^{r+s}.
\]
For a homogeneous even \(x=j/n!\), \(j\in J^n\), its rational divided power
is
\[
 \gamma_m(x)=
 \frac{(nm)!}{m!(n!)^m}\frac{j^m}{(nm)!}
       \in\frac1{(nm)!}J^{nm}.
\]
The leading coefficient is an integer: it counts partitions of an
\(nm\)-element set into \(m\) unordered blocks of size \(n\).
The elements already in \(A\) have their given divided powers, which
agree with the rational formula because \(A\) is torsion-free.
The divided-power addition formula then proves closure for every
finite sum in (136.PS1). Functoriality is immediate from preservation of
the positive-form ideal.

Every product of \(n\) positive forms has form degree at least \(n\).
This proves the finite-denominator assertion and
\(\operatorname{Sat}_{+}(A)^0=A^0\),
\(\operatorname{Sat}_{+}(A)^1=A^1\).

At simplex two, write
\[
 B_2=\Gamma_{\mathbb Z}(x_0,x_1,x_2)
                 \otimes\Lambda(e_1,e_2),
 \qquad dx_i=e_i,\quad e_0=-e_1-e_2 .
\]
There are no forms above degree two. Moreover
\((J^2)^2=B_2^2\), since every element \(f e_1e_2\) is the
product \((f e_1)e_2\). Thus
\[
 \operatorname{Sat}_{+}(B)_2^1=B_2^1,\qquad
 \operatorname{Sat}_{+}(B)_2^2=\tfrac12 B_2^2 .
\]
The standard integral PD Poincaré contraction gives
\(d(B_2^{w,1})=B_2^{w,2}\) for \(w\geq2\).
Consequently the degree-two cohomology is
\[
  \frac{\tfrac12 B_2^{w,2}}{B_2^{w,2}}
       \cong B_2^{w,2}/2 .
\]
The free abelian module \(B_2^{w,2}\) has basis
\(\gamma_{a_0}(x_0)\gamma_{a_1}(x_1)\gamma_{a_2}(x_2)e_1e_2\)
with \(a_0+a_1+a_2=w-2\), hence rank \(\binom w2\). This proves (136.PS2).
\end{proof}

\paragraph{Why this test is relevant.}
This construction permits factorial divisions of products of positive
forms while retaining strict PD closure. Its failure on \(B\) comes
from closed products of degree-one forms:
division of \(dx_1dx_2\) by two introduces a new closed two-form without
enlarging the one-form primitives. Thus finite denominators in each form
degree and strict PD closure are insufficient by themselves. The result
does not exclude a relative saturation with additional coherent
primitive choices.

